\documentclass[10pt,letterpaper]{article}
\usepackage[T1]{fontenc}
\usepackage[utf8]{inputenc}
\usepackage{lmodern,amsmath,amssymb,amsthm,mathtools}
\usepackage[margin=1in]{geometry}
\usepackage{microtype}
\usepackage[hidelinks]{hyperref}
\hypersetup{pdftitle={Global Hessian realizations of statistical manifolds},pdfauthor={Alec Kriebel}}
\newtheorem{theorem}{Theorem}
\newtheorem{lemma}[theorem]{Lemma}
\newtheorem{corollary}[theorem]{Corollary}
\newcommand{\R}{\mathbb R}
\newcommand{\dd}{\mathrm d}
\newcommand{\Hess}{\operatorname{Hess}}
\newcommand{\tr}{\operatorname{tr}}
\newcommand{\HS}{\mathrm{HS}}
\DeclareMathOperator{\Sym}{Sym}
\title{Global Hessian realizations of statistical manifolds}
\author{Alec Kriebel\\\small Independent researcher\\
\small\href{https://orcid.org/0009-0001-9320-500X}{ORCID: 0009-0001-9320-500X}}
\date{October 15, 2026}
\begin{document}
\maketitle
\begin{abstract}
Let $M$ be a smooth, Hausdorff, second-countable $n$-manifold without
boundary, with a positive Riemannian metric and two prescribed, mutually
dual, torsion-free connections. For $n\geq1$ we construct a proper smooth
embedding into an open subset of $\R^{\binom{2n+4}{3}}$ carrying one smooth
positive Hessian potential. Orthogonal projection of the ambient standard
connection and its metric dual induces exactly the ordered pair of given
connections. Finite polynomial features provide a smooth global scalar
third-jet lifting; a variable tubular neighborhood and a pointwise Schur
estimate provide ambient positivity without compactness or uniform bounds.
For a general Riemannian dual pair, vanishing of both torsions is necessary
and sufficient for such realization. This supplies a self-contained answer
in the smooth statistical category to the historical dually flat realization
question. Historical novelty and present openness are not asserted; important
prior constructions and unresolved primary-source coverage are specified.
\end{abstract}
\noindent\textbf{Status and disclosure.}
Unrefereed research note. AI tools were used extensively in developing,
writing, checking, and adversarially reviewing the argument and supporting
computations. These checks do not constitute independent human peer review
or formal proof certification. The author is responsible for the claims.

\section{Statement and conventions}

Amari asked whether an arbitrary manifold with a Riemannian dual structure
can be realized as a higher-dimensional submanifold of a dually flat
manifold with the induced dual structure \cite[p.~87, question~1]{Amari}.
The question appears in Furuhata--Matsuzoe--Urakawa
\cite[p.~125, item~1(a)]{FMU}, reporting the December 1996 workshop.
The cited definition requires zero torsion and curvature for both ambient
connections, with dual affine coordinates understood locally
\cite[pp.~84--85]{Amari}. No fixed probability model, convexity, completeness,
or globally injective dual coordinate map is imposed here.

Throughout, manifolds and data are smooth. A \emph{dual pair} satisfies
\begin{equation}\label{eq:duality}
 Xg(Y,Z)=g(\nabla_XY,Z)+g(Y,\nabla^*_XZ).
\end{equation}
It is \emph{statistical} when both connections are torsion-free. The
connections induced along an isometric immersion are the metric-orthogonal
tangential projections of the respective ambient connections. Flatness below
refers to these affine connections, not to Levi--Civita curvature.

\begin{theorem}\label{thm:main}
Let $(M,g,\nabla,\nabla^*)$ be a statistical manifold of fixed finite
dimension $n\geq1$, Hausdorff and second-countable, without boundary, with
$g$ positive definite. There exist an open set $U\subset\R^N$,
$N=\binom{2n+4}{3}>n$, a smooth function $\Psi:U\to\R$ with positive
ordinary Hessian $G=\Hess\Psi$, and a proper smooth embedding $f:M\to U$
such that $f^*G=g$. The standard flat connection $D$ and its $G$-dual $D^*$
are torsion-free and flat and induce, in that order, $\nabla$ and $\nabla^*$.
\end{theorem}

There is no connectedness or completeness hypothesis. The domain $U$ depends
on all the data and may be nonconvex and incomplete. The theorem neither
constructs a finite probability model nor optimizes $N$. Its proof uses the
proper Whitney embedding theorem and the smooth Euclidean tubular
neighborhood theorem; the jet lifting and positivity argument are given
explicitly below.

\section{A global pair from scalar third jets}

Let $L$ be the Levi--Civita connection and set $C=\nabla g$. Duality and
the two torsion conditions imply that $C$ is totally symmetric and
\begin{equation}\label{eq:difference}
 \nabla=L-\tfrac12 C^\sharp,\qquad
 \nabla^*=L+\tfrac12 C^\sharp,
 \quad g(C^\sharp(X,Y),Z)=C(X,Y,Z).
\end{equation}
Indeed, write $K=\nabla-L$ and $K^*=\nabla^*-L$. Both are symmetric in
their two vector arguments, and duality says
$g(K_XY,Z)=-g(Y,K^*_XZ)$. Applying these symmetries successively makes
$g(K_XY,Z)$ symmetric in all three variables; it follows that
$K^*=-K$ and $C(X,Y,Z)=-2g(K_XY,Z)$.

In coordinates write $g_{ij}=g(\partial_i,\partial_j)$ and
$\Gamma_{ijk}=g(\nabla_{\partial_i}\partial_j,\partial_k)$.
The symmetric array
\begin{equation}\label{eq:Q}
 Q_{ijk}:=\partial_k g_{ij}+\Gamma_{ijk}
 =\tfrac12\bigl(\partial_i g_{jk}+\partial_j g_{ik}
                         +\partial_k g_{ij}-C_{ijk}\bigr)
\end{equation}
will prescribe \emph{ordinary} third derivatives, not a covariant tensor.

\subsection*{Finite features and a smooth right inverse}

Choose a proper smooth embedding $e:M\to\R^K$, $K=2n+1$. Whitney's
proper version is recorded in \cite[Lemma~19, p.~665 and footnote~32]{Whitney};
see also its Theorem~1, p.~654. Let $p_1,\ldots,p_N$ be all monomials of
degree at most three in $K$ variables, including the constant, and put
\begin{equation}\label{eq:feature}
 f=(p_1\circ e,\ldots,p_N\circ e),\qquad
 N=\dim\mathcal P_3(\R^K)=\binom{K+3}{3}.
\end{equation}
Projection $\pi$ onto the degree-one coordinates gives $\pi f=e$.
Consequently $f$ is an injective immersion and a topological embedding.
For compact $B\subset\R^N$, $f^{-1}(B)$ is a closed subset of the compact
set $e^{-1}(\pi B)$, so $f$ is proper. Thus $S=f(M)$ is a closed embedded
submanifold of $\R^N$.

For the vector bundle of scalar third jets define the smooth bundle map
\[
 E:M\times\R^N\longrightarrow J^3(M,\R),\qquad
 E_p(c)=j^3_p\Bigl(\sum_A c_A f^A\Bigr).
\]
Each $E_p$ is onto. Choose $n$ ambient linear functions whose restrictions
by $e$ have independent differentials at $p$. These restrictions are local
coordinates $y$ near $p$. Every scalar third jet is a polynomial of degree
at most three in $y-y(p)$, hence is the restriction of an ambient polynomial
of that degree. This proves surjectivity using the same finite polynomial
space at every point, even though the coordinates used to check it vary.

Choose smooth positive bundle metrics, available by paracompactness, and
denote the adjoint by $E^\dagger$. Then
\begin{equation}\label{eq:rightinverse}
 R=E^\dagger(EE^\dagger)^{-1}
\end{equation}
is a smooth right inverse: $EE^\dagger$ is positive definite on each jet
fiber. Inversion is smooth on invertible matrices; no uniform bound for
$R$, finite chart cover, or compactness is needed.

For example, with $M=(0,\infty)$ and $e(t)=(\log t,0,0)$, the polynomial
\[
 q_t(y)=\frac{t^3}{6}(y-\log t)^3
\]
restricts to $q_t(\log s)$ with zero derivatives of orders zero through two
at $s=t$ and third derivative one. Its coefficients are smooth in $t$,
and its leading coefficient is unbounded. This illustrates a smooth lift
without a uniform bound; it is not claimed to be the minimum-norm lift
in \eqref{eq:rightinverse}.

\subsection*{The prescribed jet and its coordinate law}

At $p$, use the exponential map of $L$ only as a germ near $0\in T_pM$,
and prescribe the scalar third jet $J_p$ by
\begin{equation}\label{eq:germ}
 q_p(\exp_p v)=\tfrac12 g_p(v,v)-\tfrac1{12}C_p(v,v,v).
\end{equation}
The contractions are basis independent, and the local exponential map is
smooth near the zero section. This defines a smooth section $J$; no common
injectivity radius or complete exponential map is required. It is a section
of jets of germs at varying points, not the third jet of one global scalar
function on $M$. In particular, there is no implication from $q_p(p)=0$
for all $p$ that these germs have vanishing Hessians.

In $L$-normal coordinates the first derivative is zero, the Hessian is $g$,
and the third derivative is $-C/2$. In arbitrary coordinates the conversion
adds
\[
 L_{ij}^{a}g_{ak}+L_{ik}^{a}g_{aj}+L_{jk}^{a}g_{ai}
 =\tfrac12(\partial_i g_{jk}+\partial_j g_{ik}+\partial_k g_{ij}).
\]
Thus the value, first, second, and third ordinary derivatives of $J_p$ are
respectively $0,0,g,Q$. Explicitly, under $x=x(y)$, set
$X_a^i=\partial x^i/\partial y^a$ and
$X_{ab}^i=\partial^2x^i/(\partial y^a\partial y^b)$. The chain rule gives
\begin{equation}\label{eq:chartlaw}
 Q'_{abc}=X_a^iX_b^jX_c^k Q_{ijk}
 +g_{ij}\bigl(X_{ab}^iX_c^j+X_{ac}^iX_b^j+X_{bc}^iX_a^j\bigr).
\end{equation}
There is no third coordinate-change term because the first derivative is
zero. The connection transformation supplies the first of these three
inhomogeneous terms, and differentiating the transformed metric supplies
the other two. This checks directly why \eqref{eq:Q} is compatible scalar
jet data and must not be treated as a tensor.

Set $a=RJ$ and $\phi=-a$, identifying coefficient vectors with ambient
covectors. By definition of $E$, in every coordinate chart
\begin{equation}\label{eq:jetidentities}
 a_A f^A=0,\qquad a_A f_i^A=0,\qquad
 a_A f_{ij}^A=g_{ij},\qquad a_A f_{ijk}^A=Q_{ijk}.
\end{equation}
Here $f_i=\partial_i f$, etc. In evaluating the jet at $p$, the coefficient
$a(p)$ is held fixed in the germ variable. Subsequent differentiation along
the base $p$ differentiates $a$ as well. Differentiating the second and
third identities and using \eqref{eq:Q} gives, with the signs displayed,
\begin{align}
 f_i^A\partial_j\phi_A
   &=a_A f_{ij}^A=g_{ij},\label{eq:pairmetric}\\
 f_{ij}^A\partial_k\phi_A
   &=a_A f_{ijk}^A-\partial_k g_{ij}=\Gamma_{ijk}.
   \label{eq:pairconnection}
\end{align}
Moreover
\begin{equation}\label{eq:zeroform}
 \phi_A\dd f^A=0.
\end{equation}
These are the metric and connection identities of a Lauritzen-type pair,
with an identically zero compatibility one-form, rather than a merely
closed form whose periods would need to be controlled. The metric identity
makes $\dd\phi$ injective on each tangent space, but global injectivity of
$\phi$ is neither inferred nor needed.

\section{One positive Hessian neighborhood of the whole image}

Identify $M$ with $S$. By the smooth Euclidean normal tubular neighborhood
theorem there is an open neighborhood $T$ of $S$ on which normal addition
is a diffeomorphism from a neighborhood of the zero section of $NS$.
Write
\[
 z=r(z)+\nu(z),\qquad r(z)\in S,\quad \nu(z)\in N_{r(z)}S.
\]
We use the variable-radius version, not a uniform-radius assertion; an
explicit statement is \cite[Definition~3.53 and Theorem~3.54, pp.~43--44]{Tube}.
At $s\in S$, $\dd r=P_{T_sS}$ and $\dd\nu=P_{N_sS}$, as follows from
the derivative of normal addition at $(s,0)$. No normal trivialization is
assumed.

By \eqref{eq:zeroform}, $\phi_s$ annihilates $T_sS$. Define on $T$
\begin{equation}\label{eq:psi0}
 \Psi_0(z)=\phi_{r(z)}(\nu(z)).
\end{equation}
It is smooth, $\Psi_0|_S=0$, and $\dd\Psi_0|_S=\phi$: the term
containing the derivative of $\phi$ vanishes at $\nu=0$, while
$\phi_s P_{N_sS}=\phi_s$. Differentiating this ambient covector equality
along $S$ gives
\begin{equation}\label{eq:hessian-column}
 H_0(f_i,\,\cdot\,)=\partial_i\phi,
 \qquad H_0=\Hess\Psi_0.
\end{equation}
The tangent block is $g$ by \eqref{eq:pairmetric}. On a fixed normal fiber
\eqref{eq:psi0} is linear, so the normal--normal block is zero. In Euclidean
orthonormal tangent and normal frames,
\[
 H_0|_S=\begin{pmatrix}A&B\\B^{\mathsf T}&0\end{pmatrix},\qquad A>0.
\]
Regard $A$ as a self-adjoint endomorphism of $TS$, and $B:NS\to TS$.
The functions
\begin{equation}\label{eq:lambda}
 t=\tr(A^{-1})\|B\|_{\HS}^2,\qquad \lambda=1+t
\end{equation}
are smooth and frame independent. They need only be finite pointwise;
they may be unbounded on $S$.

Set
\begin{equation}\label{eq:psi}
 \Psi(z)=\Psi_0(z)+\lambda(r(z))|\nu(z)|^2.
\end{equation}
The correction has zero value and first derivative on $S$. Its Hessian
there is $2\lambda\langle P_N\,\cdot\,,P_N\,\cdot\,\rangle$, so
\begin{equation}\label{eq:block}
 \Hess\Psi|_S=\begin{pmatrix}A&B\\B^{\mathsf T}&2\lambda I\end{pmatrix}.
\end{equation}
For $w\in NS$, positivity of $A^{-1}$ implies
\[
 \langle B^{\mathsf T}A^{-1}Bw,w\rangle
 \leq \|A^{-1}\|_{\rm op}\|B\|_{\HS}^2|w|^2
 \leq t|w|^2.
\]
For a tangent vector $v$ the quadratic form of \eqref{eq:block} is
\begin{align*}
 &\langle A(v+A^{-1}Bw),v+A^{-1}Bw\rangle
       +\langle(2\lambda I-B^{\mathsf T}A^{-1}B)w,w\rangle,\\[-2pt]
 &\hspace{35mm}2\lambda I-B^{\mathsf T}A^{-1}B\ \geq\ (2+t)I>0.
\end{align*}
It is therefore positive definite at every point of $S$. Define
\begin{equation}\label{eq:U}
 U=\{z\in T:(\Hess\Psi)_z\text{ is positive definite}\}.
\end{equation}
This is open and contains all of $S$. Restricting $\Psi$ to $U$ produces
one global positive Hessian metric. No uniform tubular radius, lower
eigenvalue bound, bounded coefficient field, or infinite-series convergence
is used. The map $f:M\to U$ remains proper, since a compact subset of $U$
is also compact in $\R^N$.

\section{Flatness, the ordered pair, and the torsion criterion}

In ambient coordinates $z^A$, let $D$ be the standard connection and
$G_{AB}=\partial_A\partial_B\Psi$. Its metric dual has coefficients
\begin{equation}\label{eq:dualcoeff}
 (\Gamma^*)^C_{AB}=G^{CE}\partial_A G_{BE}.
\end{equation}
Symmetry of the third derivatives makes $D^*$ torsion-free. The functions
$\eta_E=\partial_E\Psi$ have invertible Jacobian $G$ and satisfy
\[
 \partial_A\partial_B\eta_E
       -(\Gamma^*)^C_{AB}\partial_C\eta_E=0.
\]
They are therefore local $D^*$-affine coordinates. This proves that $D^*$
is flat everywhere, without assuming that $z\mapsto\dd\Psi(z)$ is globally
injective. The connection $D$ is already torsion-free and flat.

The preserved first jet on $S$ gives
\begin{equation}\label{eq:gradient}
 \dd\Psi(f(p))=\phi(p),\qquad
 G_{AB}(f(p)) f_k^B=\partial_k\phi_A.
\end{equation}
Combining \eqref{eq:gradient} with \eqref{eq:pairmetric} proves $f^*G=g$.
For the connection $\nabla^{\rm ind}$ obtained from $D$,
\[
 g(\nabla^{\rm ind}_{\partial_i}\partial_j,\partial_k)
 =G(f_{ij},f_k)=f_{ij}^A\partial_k\phi_A=\Gamma_{ijk}.
\]
Thus $\nabla^{\rm ind}=\nabla$. Restricting ambient duality to tangent
fields, normal terms disappear by orthogonality. The induced connection
from $D^*$ is consequently the $g$-dual of $\nabla$, hence is the specified
$\nabla^*$. This proves Theorem~\ref{thm:main}, including the ordering.

\begin{corollary}\label{cor:torsion}
For a smooth positive Riemannian dual pair on a manifold as above,
realization by orthogonal projection from a torsion-free dually flat
statistical manifold is possible if and only if both given torsions vanish.
\end{corollary}
\begin{proof}
Sufficiency is the theorem. Along any immersion into an ambient manifold,
the torsion-free identity
$\bar\nabla_X\dd f(Y)-\bar\nabla_Y\dd f(X)=\dd f([X,Y])$
projects to the torsion-free identity for the induced connection.
Apply it separately to the two ambient connections.
\end{proof}

A second-countable zero-dimensional manifold is countable and discrete.
Embed it as a subset of the nonnegative integers in $\R$, with
$\Psi(z)=z^2/2$. This is proper and realizes all its (zero-dimensional)
tangent data. The empty case is vacuous. No extension to boundary or
indefinite-metric settings is asserted.

\section{Prior constructions and the scope of the answer}

The conclusion answers the historical existence question in its smooth
positive statistical category. It is not a claim that this is the first
answer or that the question remained open until this note. Several
ingredients have substantial precedents. Lauritzen's statistical,
cubic-tensor, dual-connection, and induced-connection framework is in
\cite{LauritzenIMS}. The separate Lancaster note \cite{LauritzenPair}
is credited for pair embeddings by Marugame
\cite[\S2.4]{Marugame}; its original body was not obtained in the present
source audit. Matumoto \cite[Theorem~1 and proof, pp.~328--331]{Matumoto}
constructs global contrast functions by compatible scalar third-order data
and explicitly distinguishes the further flat-embedding question.

Finite jets, freedom, and right inverses belong to the Nash--Gromov
embedding tradition. In statistical geometry they occur in L\^e's long
preprint \cite[pp.~33--39]{Le2003} and in Loi--Matta's robustness theorem
and minimum-norm inverse \cite{LoiMatta}. The latter starts with already
induced statistical data; openness near such data does not alone realize
an arbitrary prescribed metric and cubic. Marugame's local equivalence
\cite[Theorem~2.2, \S2.4]{Marugame} constructs a Hessian extension from a
pair using a local primitive and a positive transverse quadratic term.
Those are direct antecedents of the extension stage here. The present
proof supplies a global pair with $\phi\,\dd f=0$ and a variable-neighborhood
positivity argument. Osipov's positive Hessian cones are genuine global
realizations for the stated constant-statistical-curvature classes
\cite{Osipov}; their restriction is the curvature hypothesis, not the
absence of a single global affine chart.

The finite-probability-model literature needs a separate distinction.
L\^e's earlier texts make unrestricted finite-simplex claims
\cite{Le2003,Le2006,Le2007}; the inspected arXiv v6 explicitly assumes
compactness in its Main Theorem and Theorems~5.1 and~5.6 \cite{Le2016}.
In the actual 2017 book, Theorem~4.10 on p.~222 gives finite-sample
immersion for compact input and a countable-sample model for noncompact
input; Theorem~4.11 on p.~238 makes the corresponding embedding distinction
\cite{AJLS}. These statements were inspected in a Springer-authorized
preview; they should not be replaced by the broader chapter abstract.
The complete chapter, including the unavailable p.~239, was not audited.

A valid unrestricted smooth finite-simplex embedding theorem would already
answer the historical question: in mixture coordinates the potential
$\sum_a p_a\log p_a$ has positive Hessian equal to the Fisher metric, and
its derivative cubic is minus the Amari--Chentsov cubic. Applying such a
theorem to the appropriate cubic sign induces both ordered connections
by \eqref{eq:difference}. Properness is not a historical novelty shield
either: an embedded image $S$ is open in its ambient closure, so removing
the closed set $\overline S\setminus S$ makes the embedding proper in a
smaller allowed open ambient. Neither observation proves the particular
dimension bound here to be new.

The inspected older noncompact gluing arguments leave compatibility of the
separately embedded pieces and their derivatives unproved; this is not a
counterexample to their conclusions. The final 2006 article and final 2007
chapter were not obtained as authenticated full bodies, and are not assumed
identical to an arXiv version. Along with the separate Lauritzen note and
Abe's article \cite{Abe}, they remain material primary-source gaps. The
accepted manuscripts inspected for Marugame and Loi--Matta are distinguished
from their publisher versions in the references. This bounded literature
audit establishes neither absence of an earlier full solution nor current
openness. The note presents the explicit proof above with those limitations.

\paragraph{Verification materials.}
Companion materials preserve the original proof and its computational
records, independent checks of jet and Hessian identities, source-scope
comparisons, and reproduction instructions. Finite exact checks and
numerical examples corroborate identities; they do not replace the universal
proof or its differential-topology inputs. Historical strict floating-output
comparison failures are preserved separately from a typed, bounded portable
comparison of designated sample fields. Its tolerances are a transport
policy, not a proved roundoff estimate. No numerical test, AI review, or
source search certifies historical novelty.

\begin{thebibliography}{99}\small
\bibitem{Amari}
S.-I. Amari, \emph{Information geometry}, Contemporary Mathematics
\textbf{203} (1997), 81--95, especially pp.~84--85, 87.
\url{https://bsi-ni.brain.riken.jp/database/file/166/169.pdf}.

\bibitem{FMU}
H. Furuhata, H. Matsuzoe, and H. Urakawa,
\emph{Open problems in affine differential geometry and related topics},
Interdisciplinary Information Sciences \textbf{4}(2) (1998), 125--127.
\url{https://www.jstage.jst.go.jp/article/iis/4/2/4_2_125/_pdf/-char/en}.

\bibitem{LauritzenIMS}
S. L. Lauritzen, \emph{Statistical manifolds}, in
\emph{Differential Geometry in Statistical Inference}, IMS Lecture Notes
Monograph Series \textbf{10} (1987), 163--216.
\href{https://doi.org/10.1214/lnms/1215467061}{doi:10.1214/lnms/1215467061}.

\bibitem{LauritzenPair}
S. L. Lauritzen, \emph{Conjugate connections in statistical theory},
in \emph{Geometrization of Statistical Theory}, ULDM Publications,
University of Lancaster (1987), 33--51. Original body not obtained;
pair attribution checked in \cite[\S2.4]{Marugame} and the author's
\href{https://web.math.ku.dk/~lauritzen/publications.html}{institutional bibliography}.

\bibitem{Matumoto}
T. Matumoto, \emph{Any statistical manifold has a contrast function---On
the $C^3$-functions taking the minimum at the diagonal of the product
manifold}, Hiroshima Mathematical Journal \textbf{23}(2) (1993), 327--332.
\href{https://doi.org/10.32917/hmj/1206128255}{doi:10.32917/hmj/1206128255}.

\bibitem{Whitney}
H. Whitney, \emph{Differentiable manifolds}, Annals of Mathematics
\textbf{37}(3) (1936), 645--680.
\href{https://doi.org/10.2307/1968482}{doi:10.2307/1968482}.

\bibitem{Tube}
M. Gualtieri, \emph{MAT1300 course notes}, part 11 (2017),
Definition~3.53 and Theorem~3.54, pp.~43--44.
\url{https://math.toronto.edu/mgualt/courses/17-1300/docs/17-1300-notes-11.pdf}.

\bibitem{Le2003}
H. V. L\^e, \emph{Monotone invariants and embeddings of statistical
manifolds}, MPI preprint 43/2003;
\url{https://files-www.mis.mpg.de/mpi-typo3/preprints/2003/preprint2003_43.pdf}.
See also \emph{Statistical manifolds are statistical models}, MPI preprint
77/2003;
\url{https://files-www.mis.mpg.de/mpi-typo3/preprints/2003/preprint2003_77.pdf}.

\bibitem{Le2016}
H. V. L\^e, \emph{Monotone invariants and embeddings of statistical
manifolds}, arXiv:math/0506163v6, March 10, 2016, especially pp.~14--15,
22--23. \url{https://arxiv.org/abs/math/0506163v6}.
This inspected version is distinguished from \cite{Le2007}.

\bibitem{Le2006}
H. V. L\^e, \emph{Statistical manifolds are statistical models},
Journal of Geometry \textbf{84} (2006), 83--93.
\href{https://doi.org/10.1007/s00022-005-0030-0}{doi:10.1007/s00022-005-0030-0}.
Publication metadata and broad abstract inspected; final body not obtained.

\bibitem{Le2007}
H. V. L\^e, \emph{Monotone invariants and embeddings of statistical
manifolds}, in \emph{Advances in Deterministic and Stochastic Analysis},
World Scientific (2007).
\href{https://doi.org/10.1142/9789812770493_0012}{doi:10.1142/9789812770493\_0012}.
Final chapter body not obtained; no identity with an arXiv version assumed.

\bibitem{AJLS}
N. Ay, J. Jost, H. V. L\^e, and L. Schwachh\"ofer,
\emph{Information Geometry}, Springer (2017), Chapter~4,
\emph{The intrinsic geometry of statistical models}, 185--239.
\href{https://doi.org/10.1007/978-3-319-56478-4_4}{doi:10.1007/978-3-319-56478-4\_4}.
Actual Theorems~4.10 and~4.11 inspected in the authorized previews
\href{https://books.google.com/books?id=pLsyDwAAQBAJ&pg=PA222}{p.~222}
and \href{https://books.google.com/books?id=pLsyDwAAQBAJ&pg=PA238}{p.~238};
full chapter not obtained.

\bibitem{Marugame}
T. Marugame, \emph{The Bonnet theorem for statistical manifolds},
Information Geometry \textbf{4} (2021), 363--376.
\href{https://doi.org/10.1007/s41884-021-00056-4}{doi:10.1007/s41884-021-00056-4}.
Inspected bodies: \href{https://arxiv.org/abs/2103.10102v1}{arXiv:2103.10102v1}
(March 18, 2021), \S\S2.1--2.4, and the
\href{https://uec.repo.nii.ac.jp/record/10333/files/Bonnet_thm_statistical_mfd.pdf}{author accepted manuscript};
publisher version not obtained.

\bibitem{LoiMatta}
A. Loi and S. Matta, \emph{Robustness of statistical manifolds},
Topology and its Applications \textbf{327} (2023), 108438.
\href{https://doi.org/10.1016/j.topol.2023.108438}{doi:10.1016/j.topol.2023.108438}.
Inspected bodies: \href{https://arxiv.org/abs/2110.06145v1}{arXiv:2110.06145v1}
and the \href{https://iris.unica.it/retrieve/6dbf3c46-1ccf-403f-bb6f-a3757afe194a/TOPOL\%20108438_merged.pdf}{author accepted manuscript},
Theorem~1 and the free-map/right-inverse argument; publisher version not obtained.

\bibitem{Osipov}
P. Osipov, \emph{Statistical Lie algebras of a constant curvature and
locally conformally K\"ahler Lie algebras},
\href{https://arxiv.org/abs/2112.06686v2}{arXiv:2112.06686v2}
(May 20, 2022), Theorem~2.12;
\emph{Locally conformally Hessian and statistical manifolds},
\href{https://arxiv.org/abs/2209.02357v3}{arXiv:2209.02357v3}
(September 6, 2023), Theorem~3.16, p.~14.
The latter was published in Journal of Geometry and Physics,
\href{https://doi.org/10.1016/j.geomphys.2023.104989}{doi:10.1016/j.geomphys.2023.104989};
the cited theorem reading is of the specified preprint.

\bibitem{Abe}
N. Abe, \emph{Affine immersions and conjugate connections}, Tensor,
New Series \textbf{55}(3) (1994), 276--280.
\href{https://dl.ndl.go.jp/en/pid/11196378}{Issue archive}; original body not obtained.
\end{thebibliography}
\end{document}
